% Reserved candidate D.2.54. Not registered before root's independent audit.
\paragraph{Theorem D.2.54 (Proved: the simply connected two-fixed-component star-log families).}
Take one of the following six default natural rows of D.2.49, with
all other caps carrying their minimal natural marking and with no extra
regular logarithmic fibres:
\[
\begin{array}{c|c|c|c|c}
\text{row}&a&b_0&\text{zero index}&\text{star type }I_n^*\\\hline
\mathrm{OS46/default}&1&4&1,2,3&I_3^*\\
\mathrm{OS50/default}&1&4&1&I_1^*\\
\mathrm{OS50/default}&3&4&2,3&I_1^*.
\end{array}                                                   \tag{D.SS1}
\]
At the star choose any integer $h\ge1$, put $k=2h$, and take the
actual root order $m=2k=4h$. Allow every even integer component phase
$j$, every integer transverse twist $r$, and every ordinary clutch
$e$. Choose either sign $s=\pm1$ and set
\[
                         N=4he+s.                            \tag{D.SS2}
\]
These are free holomorphic affine actions by D.2.49. Their compact
outputs have the oriented diffeomorphism type
\[
 X\ \cong\ \mathop{\#}_{nh}(S^2\times S^4)
                    \ \#\ \mathop{\#}_{nh}(S^3\times S^3).
                                                               \tag{D.SS3}
\]
The assertion includes every displayed integer phase lift and twist;
it is not a claim based only on a finite range of root orders.
In particular, the six rows give two unbounded families, indexed by
$nh=h$ and $nh=3h$, of explicit complex structures on these smooth
manifolds.

More precisely, let $S_0,\ldots,S_t$ be the components of the reduced
star fibre in their interval order, where $t=nh$. Their divisor
classes $x_i$ have the exact primitive lattice and cubic
\[
 H^2(X;\mathbb Z)=
    \mathbb Z\langle x_0,\ldots,x_t\rangle/
                         \mathbb Z(x_0+\cdots+x_t),
 \qquad x_i x_j x_l=0.                                     \tag{D.SS4}
\]
Thus any $t$ of these component classes form a primitive basis.
All characteristic classes in the following display vanish integrally:
\[
             c_1(X)=c_2(X)=p_1(X)=0.                       \tag{D.SS5}
\]

\begin{proof}
\emph{The actual central components and their multiplicities.}
The canonical root has $q=nk=2nh$ semistable components. Because $q$
and $j$ are even, the reflection $i\mapsto j-i$ has two fixed
vertices and no fixed edges. No pair of adjacent vertices belongs to
one vertex orbit: that would require $2i=j\pm1\pmod q$.
Consequently the quotient reduced fibre is a chain of $t+1$
\emph{embedded} smooth ruled surfaces, with no self-intersections of
one component and no triple intersections. This includes $q=2$,
when the two fixed components meet along one conductor curve.

The two endpoints are the odd oriented $S^2$ bundles over a torus
from D.2.49. Choose their section and fibre classes $\sigma,F$ with
\[
 \sigma^2=1,\quad \sigma F=1,\quad F^2=0,\qquad
 C=2\sigma-F,\quad C^2=0,\quad K=-2\sigma+F=-C.
                                                               \tag{D.SS6}
\]
Here $C$ is the actual connected conductor bisection. Each interior
surface is a ruled surface over a torus with two disjoint pole
sections, both of class $\sigma$, where
\[
       \sigma^2=0,\qquad K=-2\sigma.                     \tag{D.SS7}
\]
These intersection and canonical-class formulas also follow from
adjunction on the fibre and on a section: $KF=-2$ and
$(K+\sigma)\sigma=0$. They are unaffected by flat unit factors.

It is essential to retain the original base coordinate. On the
semistable cover it is $z=u^{2k}$. The total quotient is etale,
because its affine action is free on every stratum. Near a generic
point of \emph{any} quotient component, a local sheet of this
etale cover identifies $u$ with a normal coordinate. Therefore
the vanishing multiplicity of $z$ is $2k$ on both fixed and paired
components, and
\[
                       \operatorname{div}(z)=2k\sum_{i=0}^t S_i.
                                                               \tag{D.SS8}
\]
There is no stabilizer-order ratio in this formula. Restrict it to
$S_i$. Since $H^2(S_i;\mathbb Z)$ is torsion-free, it gives the
actual normal class
\[
 c_1(\nu_{S_i/X})=-\sum_{j\text{ adjacent to }i}[C_{ij}]
                      =K_{S_i}.                          \tag{D.SS9}
\]
At an endpoint the sum is its bisection, and at an interior vertex
it is its two pole sections. In particular
$c_1(TX)|_{S_i}=c_1(TS_i)+c_1(\nu_{S_i/X})=0$.
This proves the needed cancellation from the actual multiplicities,
without assuming a trivial canonical bundle on the whole output.

\emph{The primitive lattice, not only its rational span.}
D.2.49 gives $\pi_1(X)=1$, $H_2=H_4=\mathbb Z^t$ and
$H_3=\mathbb Z^{2t}$ for (D.SS2). In particular $H^2(X)$ is free
of rank $t$. Formula (D.SS8) gives $\sum_i x_i=0$ integrally.
For a divisor $D=\sum u_iS_i$, its component restrictions are
\[
\begin{aligned}
 D|_{S_0}&=(u_1-u_0)(2\sigma-F),\\
 D|_{S_i}&=(u_{i-1}-2u_i+u_{i+1})\sigma
                                  &&(0<i<t),\\
 D|_{S_t}&=(u_{t-1}-u_t)(2\sigma-F).
\end{aligned}                                               \tag{D.SS10}
\]
When $t=1$ only the first and last rows occur. These formulas use
the actual Cartier branches, so each conductor is counted once.
At each endpoint select the coefficient of $-F$; at each interior
vertex select the coefficient of $\sigma$. This gives, up to signs,
the ordinary path-graph Laplacian. Every reduced cofactor of that
matrix is a unit: successively subtracting the endpoint row proves
this by induction along the path. Thus after deleting any one
component column, the full restriction matrix contains a unimodular
$t$ by $t$ minor.

It follows first that the component subgroup has rank $t$ and only
the stated all-one relation. It follows also that its image in
$\bigoplus_iH^2(S_i;\mathbb Z)$ is primitive. The restriction map
on $H^2(X)$ is rationally injective, because the component subgroup
already has full rank. If that subgroup had a proper finite-index
extension in $H^2(X)$, its restriction would be a proper integral
extension of a primitive subgroup of this direct sum, which is
impossible. This proves (D.SS4) as an exact integral lattice.

For an independent marking check, the literal multiplicative
vanishing columns in each of the six rows (D.SS1) generate the
whole vertical $\mathbb Z^3$. Removing only the star cap therefore
leaves fundamental group $\mathbb Z\gamma$: all other meridians
are filled, and the ordinary clutch expresses the star meridian
as $e\gamma$. A small component meridian in the star cap has
collapsed word $g^{2k}v^{-1}$, whose exterior image is
$(2ke-N)\gamma=-s\gamma$. The terms in the other three fibre
directions vanish there. Thus every component meridian is a unit
in this exterior group. Relative Poincare--Lefschetz duality gives
the same primitive all-one relation, with no finite enlargement.
This check retains the prescribed cusp and affine markings.

\emph{All cubic and characteristic data.}
For each component, (D.SS6)--(D.SS9) give
$S_i^3=K_{S_i}^2=0$. If $i,j$ are adjacent then
$S_i^2S_j=K_{S_i}C_{ij}=0$, since all conductor sections and
bisections have square zero. If they are not adjacent the product
vanishes. Three distinct components have no common point. Thus
every triple in the primitive lattice (D.SS4) has value zero.
The restriction map in (D.SS10) is injective over $\mathbb Q$,
so the already proved restrictions of $c_1(X)$ imply its rational
vanishing. Freeness of $H^2(X)$ gives $c_1(X)=0$ integrally.

Every ruled surface over a torus has $c_2(S_i)=0$ and
$c_1(S_i)^2=0$. The tangent-normal exact sequence and (D.SS9)
therefore give
\[
 \langle c_2(X),[S_i]\rangle
   =c_2(S_i)-K_{S_i}^2=0,\qquad
 \langle p_1(X),[S_i]\rangle
   =p_1(S_i)+K_{S_i}^2=0.                                \tag{D.SS11}
\]
The classes $[S_i]$ span $H_4(X)$ integrally by (D.SS4) and duality.
Also $H^4(X)$ is free because both $H_3(X)$ and $H_4(X)$ are free.
The integral evaluation pairing is perfect, so (D.SS11) proves
the remaining equalities in (D.SS5).

Finally apply the oriented smooth classification of simply
connected spin six-manifolds in \cite[Theorem~1 and \S1.5]{manZhubr75}.
For free $H_2=\mathbb Z^t$ its data are the rank of $H_3$, the
integral cubic and the Pontryagin class; the refined class has no
torsion ambiguity, since its multiple by four is the ordinary
Pontryagin class. Equivalently, the full fundamental-class invariant
in $H_6(K(\mathbb Z^t,2);\mathbb Z)$ is determined by the integral
cubic. Our data are $b_3=2t$, cubic zero and Pontryagin class zero.
They equal those of the connected sum in (D.SS3), which is simply
connected and spin. The classification gives the asserted oriented
diffeomorphism, rather than merely a homology identification.
\end{proof}

The exact certificate \path{code/fourfolds_star_smooth.py} retains
the actual six vertical lattices and canonical star markings,
checks the component-meridian identity symbolically for all integer
parameters, and verifies the path restriction and its primitive
Smith form. The formula proof is uniform in $h$; the recorded finite
checks are corroborations, not an upper bound on the family.
